\documentclass[11pt,a4paper]{article}
\usepackage[T1]{fontenc}
\usepackage{lmodern}
\usepackage[margin=25mm]{geometry}
\usepackage{amsmath,amssymb,amsthm,booktabs,microtype,enumitem}
\usepackage[hidelinks]{hyperref}
\hypersetup{pdftitle={A corrected modulo-sixteen classification for A396844},pdfauthor={Dev Goyal},pdfsubject={A396844; formal power series and Lean proofs}}
\setlength{\parskip}{3pt}
\setlength{\emergencystretch}{3em}
\newtheorem{theorem}{Theorem}[section]
\newtheorem{lemma}[theorem]{Lemma}
\newtheorem{corollary}[theorem]{Corollary}
\theoremstyle{remark}\newtheorem{remark}[theorem]{Remark}
\newcommand{\Z}{\mathbb Z}
\newcommand{\N}{\mathbb N}
\newcommand{\F}{\mathbb F}
\newcommand{\coeff}[2]{[x^{#1}]#2}
\title{A corrected modulo-sixteen classification for A396844}
\author{Dev Goyal}
\date{24 October 2026}
\begin{document}
\maketitle
\begin{center}\fbox{\begin{minipage}{0.91\linewidth}
\textbf{AI disclosure and review status.}
ChatGPT (OpenAI) generated the mathematical arguments, formal proof code,
and this manuscript from Dev Goyal's prompts. Dev Goyal is the sole human
author and initiated the investigation. This consolidated presentation is
AI-generated; independent human review of this presentation is pending.
Source conjectures and prior results retain the credits given below.
\end{minipage}}\end{center}

\begin{abstract}
We give a separate complete manuscript for A396844, including the
parameter-family proof on which its conclusions depend. For the unique
normalized solution of $A_m(xA_m-mxA_m^2)=x^2$, we prove parity,
exact residues modulo $m^2$, binary-digit divisibility for both the series
and its compositional inverse, and the minimum prime-exponent criterion.
The entry under consideration is $m=4$.
\end{abstract}
\section{The entry and the claims}
Let $a_n(m)=[x^n]A_m(x)$, where
\begin{equation}\label{eq:general}
 A_m\bigl(xA_m(x)-mxA_m(x)^2\bigr)=x^2,
 \qquad A_m=x+O(x^2).
\end{equation}
The sequence A396844 is the case $m=4$; the constant
coefficient is zero, and the OEIS sequence is indexed from one.
The defining series and source questions are due to Paul D. Hanna
\cite{oeis4}.
We write $s_2(n)$ for the number of ones in the binary expansion of $n$.
The general results below give, in particular,
\[
 m^{s_2(n-1)}\mid a_n(m)\quad(n\ge1),
 \qquad
 a_{2^j+1}(m)\equiv m(-1)^j\pmod{m^2}.
\]
All other coefficients of degree at least two vanish modulo $m^2$.
The binary bound also holds for the coefficients of the inverse series.

The corrected residue theorem is, for $n\ge2$,
\[
a_n(4)\equiv
\begin{cases}
4\pmod{16},&n=4^k+1,\ k\ge0,\\
12\pmod{16},&n=2\cdot4^k+1,\ k\ge0,\\
0\pmod{16},&\text{otherwise}.
\end{cases}
\]
The first exceptional class is essential. The earlier source assertion
specified residue twelve in the second class and divisibility by sixteen
otherwise for $n>8$. At $n=17=4^2+1$ the correct residue is four,
so that other clause is false. We prove the corrected theorem, not the
false original clause. All coefficients after the first are even.

\paragraph{Provenance.}
These arguments are adapted from the original A396843 repository supplied
by Dev Goyal. Its historical declaration of human verification concerns
that earlier report. It is not a declaration that this new editorial
consolidation has received independent human review.
The parameter-family results are retained mathematical work, not new
problems introduced by this final packaging.
\section{Formal power series and notation}\label{sec:prelim}

\subsection{Operations and coefficient comparison}

A formal power series is an expression
$U(x)=\sum_{n\ge0}u_nx^n$ whose coefficients are considered algebraically.
We do not substitute numerical values of $x$, and no convergence assumption
is needed. Equality means equality of every coefficient. We write
$[x^n]U=u_n$. Addition is coefficientwise, and multiplication is defined by
\[
 [x^n](UV)=\sum_{i=0}^{n}u_i v_{n-i}.
\]
Thus each coefficient of a product is a finite sum. An empty sum is zero.

If $V(0)=0$, composition is defined by
\[
 U(V(x))=\sum_{n\ge0}u_n V(x)^n.
\]
The power $V(x)^n$ has no terms of degree below $n$, so each coefficient
again receives only finitely many contributions. For example, if
$U(x)=x+x^2$ and $V(x)=x^2$, then $U(V(x))=x^2+x^4$.
Composition of series with zero constant coefficient is associative:
one can verify this in each degree by the usual rules for polynomials,
since only finitely many terms contribute.

The notation $O(x^d)$ denotes an unspecified series whose coefficients below
degree $d$ vanish. In particular, $U=x+O(x^2)$ means that $U$ has constant
coefficient zero and coefficient of $x$ equal to one. We call this
\emph{normalized}. Products and compositions satisfy
\[
 O(x^a)O(x^b)=O(x^{a+b}),\qquad
 U(V(x))=O(x^{ab})
 \quad\text{if }U=O(x^a),\ V=O(x^b),\ b\ge1.
\]
These statements specify vanishing of coefficients, not estimates as a
numerical variable tends to zero.

We occasionally allow coefficients in a commutative ring $R$ with identity:
a coefficient system in which addition, subtraction, and multiplication
obey their usual algebraic laws, with a multiplicative identity $1$.
The examples used here are the integers $\Z$, residues $\Z/q\Z$, and
integer polynomials $\Z[t]$. The notation $R[[x]]$ denotes formal power
series with coefficients in $R$.

For integer series, $U\equiv V\pmod q$ means that every coefficient of
$U-V$ is divisible by $q$. Reducing coefficients modulo $q$ commutes with
addition, multiplication, and every composition used here: each coefficient
is obtained by a finite sequence of integer operations. In particular,
$\F_2=\Z/2\Z$ has the two elements $0$ and $1$, and addition equals
subtraction in this ring.

If $U(0)=0$, the series $1+U$ has a multiplicative inverse, given formally by
\[
 (1+U)^{-1}=\sum_{k\ge0}(-U)^k.
\]
The expression is well defined coefficientwise, and multiplication makes
successive terms cancel. This justifies fractions such as $x/(1+x)$.

\subsection{Compositional inverses}

A compositional inverse reverses substitution rather than multiplication.
We will need both identities in the following elementary fact.

\begin{lemma}\label{lem:inverse}
For any commutative ring $R$ with identity and any normalized series
$A\in R[[x]]$, there is a unique normalized series $B\in R[[x]]$ such that
\[
 A(B(x))=B(A(x))=x.
\]
\end{lemma}

\begin{proof}
Write $A=x+\sum_{n\ge2}a_nx^n$ and seek
$B=x+\sum_{n\ge2}b_nx^n$ with $A(B)=x$.
In degree $n\ge2$, the linear term $B$ contributes $b_n$.
In every higher power $B^k$, with $k\ge2$, a contribution to degree $n$
uses only coefficients of $B$ of index at most $n-k+1<n$.
Consequently, the coefficient equation in degree $n$ determines $b_n$
uniquely from earlier coefficients, using only addition, subtraction, and multiplication.
This constructs the unique normalized $B$ with $A\circ B=x$.

Apply the same construction to $B$ to obtain a normalized $C$ with
$B\circ C=x$. Associativity gives
\[
 A=A\circ(B\circ C)=(A\circ B)\circ C=C.
\]
Hence $B\circ A=B\circ C=x$ as well. Uniqueness follows from the
construction of the right inverse.
\end{proof}

\section{Existence and uniqueness of the defining series}

The following coefficient comparison is the reason modular solutions can
be identified by direct substitution.

\begin{lemma}\label{lem:unique}
Let $R$ be a commutative ring with identity and let $m\in R$. There is exactly
one normalized series $A\in R[[x]]$ satisfying
\[
 A(xA-mxA^2)=x^2.
\]
\end{lemma}

\begin{proof}
For a normalized series $A$, set
$T_A=xA-mxA^2$ and $\Psi_m(A)=A(T_A)$.
We first show how a change in one coefficient affects the defining equation:
\begin{equation}\label{eq:perturb}
 B-A=cx^r+O(x^{r+1}),\quad r\ge2
 \quad\Longrightarrow\quad
 \Psi_m(B)-\Psi_m(A)=cx^{r+1}+O(x^{r+2}).
\end{equation}
Indeed,
\[
 T_B-T_A=x(B-A)\bigl(1-m(A+B)\bigr)
 =cx^{r+1}+O(x^{r+2}).
\]
Both $T_A$ and $T_B$ start with $x^2$. If
$A=x+\sum_{k\ge2}a_kx^k$, then
\[
 A(T_B)-A(T_A)=(T_B-T_A)+\sum_{k\ge2}a_k(T_B^k-T_A^k).
\]
For $k\ge2$, factor $T_B^k-T_A^k$ as $T_B-T_A$ times a sum of
products of $k-1$ inner arguments. Its degree is therefore at least
$r+1+2(k-1)\ge r+3$. It follows that
\[
 A(T_B)-A(T_A)=(T_B-T_A)+O(x^{r+3}).
\]
The other difference, $(B-A)(T_B)$, is $O(x^{2r})$, hence $O(x^{r+2})$.
Adding the two differences proves \eqref{eq:perturb}.

For uniqueness, if two solutions were different, take the smallest index
$r$ at which their coefficients differ. In \eqref{eq:perturb}, $c$ is then
nonzero, so their images under $\Psi_m$ differ in degree $r+1$.
They cannot both have image $x^2$.

For existence, start with $A=x$, for which $\Psi_m(A)=x^2-mx^3$.
Suppose the residual $\Psi_m(A)-x^2$ has already been made zero through
degree $r$. Let $c$ be its coefficient of $x^{r+1}$. Replacing $A$ by
$A-cx^r$ removes that coefficient without changing any earlier one, by
\eqref{eq:perturb}. Continue for $r=2,3,\ldots$. Each coefficient of $A$
eventually stops changing. Each coefficient of its defining equation depends
on only finitely many coefficients of $A$, so the resulting series satisfies
the equation in every degree.

More precisely, for $d\ge2$, the coefficient $[x^d]\Psi_m(A)$ depends
only on $a_0,\ldots,a_{d-1}$. The inner series $T_A$ vanishes below
degree two, so only outer indices $k\le d/2$ contribute. Each coefficient
of $T_A$ through degree $d$ depends only on coefficients of $A$ through
degree $d-1$. Thus later corrections to $A$ cannot change a coefficient
of the defining equation that has already been fixed.
\end{proof}

The proof uses no division by a coefficient. In particular it applies in
residue rings even when the modulus is composite. Applying the lemma over
$\Z$ defines $A_m$ for every integer $m$. Its first terms are
\begin{equation}\label{eq:first}
 A_m(x)=x+mx^2+(2m^2-m)x^3+(5m^3-2m^2)x^4+O(x^5).
\end{equation}
For example, the coefficient of $x^3$ in the defining composition is
$a_2(m)-m$, so $a_2(m)=m$.

Applying the same construction with an indeterminate parameter $t$ over
$\Z[t]$ shows that each $a_n(t)$ is an integer polynomial. Substitution
$t=m$ then gives $a_n(m)$, by uniqueness. This observation will also allow
a polynomial version of the divisibility theorem.

\section{Parity}

\begin{theorem}\label{thm:parity}
If $m$ is odd, every $a_n(m)$ with $n\ge1$ is odd. If $m$ is even, then
$a_1(m)=1$ and every $a_n(m)$ with $n\ge2$ is even.
\end{theorem}

\begin{proof}
Suppose first that $m$ is odd. Write $P$ for the series obtained by reducing
each coefficient of $A_m$ modulo two. Equation \eqref{eq:general} becomes
\begin{equation}\label{eq:char2}
 P\bigl(xP(x)+xP(x)^2\bigr)=x^2
 \quad\text{in }\F_2[[x]],\qquad P=x+O(x^2).
\end{equation}
The normalized solution is unique by Lemma~\ref{lem:unique}.
Consider the explicit series
\[
 H(x)=\frac{x}{1+x}=\sum_{n\ge1}x^n\quad\text{in }\F_2[[x]].
\]
The last equality is the geometric series identity from
Section~\ref{sec:prelim}, with all signs equal modulo two.
Direct calculation gives
\[
 xH(x)+xH(x)^2
 =\frac{x^2(1+x)+x^3}{(1+x)^2}
 =\frac{x^2}{(1+x)^2}=:T(x).
\]
Therefore
\[
 H(T(x))=\frac{T(x)}{1+T(x)}
 =\frac{x^2}{(1+x)^2+x^2}=x^2,
\]
since $(1+x)^2+x^2=1$ in $\F_2[[x]]$.
Thus $H$ satisfies \eqref{eq:char2}, and uniqueness gives $P=H$.
Every nonconstant coefficient of $H$ is one, so every $a_n(m)$ with
$n\ge1$ is odd.

If $m$ is even, the reduced equation is $P(xP(x))=x^2$.
The series $P=x$ satisfies it and is the unique normalized solution by
Lemma~\ref{lem:unique}. Its coefficients after the first are zero,
which proves the assertion for even $m$.
\end{proof}

Taking $m=3$ proves the coefficient-parity conjecture in OEIS A396843.

\section{An exact formula modulo \texorpdfstring{$m^2$}{m squared}}

The parity proof identifies a solution after reduction modulo two. We now
use the same uniqueness principle modulo the parameter $m$, and retain
one further level of information modulo $m^2$.

\begin{theorem}\label{thm:modsquare}
For every integer $m\ge2$,
\begin{equation}\label{eq:modsquare}
 A_m(x)\equiv x+m\sum_{j\ge0}(-1)^jx^{2^j+1}\pmod{m^2}.
\end{equation}
Equivalently, for $n\ge2$,
\begin{equation}\label{eq:coeffmod}
 a_n(m)\equiv
 \begin{cases}
   m(-1)^j\pmod{m^2},&n=2^j+1\text{ for some }j\ge0,\\
   0\pmod{m^2},&\text{otherwise}.
 \end{cases}
\end{equation}
\end{theorem}

\begin{proof}
Reduce \eqref{eq:general} modulo $m$. The resulting equation is
$P(xP)=x^2$, whose unique normalized solution is $P=x$.
Hence $A_m\equiv x\pmod m$, and we can write
\[
 A_m=x+mC,\qquad C\in x^2\Z[[x]].
\]
The notation $C\in x^2\Z[[x]]$ means that $C$ has integer coefficients
and no terms of degree below two. Working modulo $m^2$ gives
\[
 T:=xA_m-mxA_m^2
 \equiv x^2+m(xC-x^3).
\]
Since $T\equiv x^2\pmod m$, we also have
$C(T)\equiv C(x^2)\pmod m$. Consequently,
\[
 A_m(T)=T+mC(T)
 \equiv x^2+m\bigl(xC-x^3+C(x^2)\bigr)\pmod{m^2}.
\]
Equation \eqref{eq:general} implies
\begin{equation}\label{eq:linear}
 xC(x)+C(x^2)\equiv x^3\pmod m.
\end{equation}
Here cancellation of $m$ is ordinary integer divisibility: if $m^2$ divides
$md$, then $m$ divides $d$. It is applied separately to each coefficient.

Set
\[
 S(x)=\sum_{j\ge0}(-1)^jx^{2^j+1}
 =x^2-x^3+x^5-x^9+x^{17}-\cdots.
\]
The exponents increase, so this is a well-defined formal series. Its
successive terms give
\[
 xS(x)=x^3-x^4+x^6-x^{10}+\cdots,\qquad
 S(x^2)=x^4-x^6+x^{10}-\cdots.
\]
Thus $xS+S(x^2)=x^3$ exactly.

The solution of \eqref{eq:linear} with no terms below degree two is unique
over $\Z/m\Z$. Indeed, if a nonzero difference $D$ began in degree
$d\ge2$, then $xD$ would begin in degree $d+1$, while $D(x^2)$ would
begin in degree $2d>d+1$. Its first term could not cancel in
$xD+D(x^2)$. Therefore $C\equiv S\pmod m$, and
$A_m=x+mC$ yields \eqref{eq:modsquare}.
\end{proof}

\section{Divisibility from binary digits}

\subsection{The binary digit inequality}

Every nonnegative integer $r$ has a unique binary expansion as a sum of
distinct powers of two. Repeated division by two gives the binary digits as
successive remainders, which also makes the expansion unique.
Define $s_2(r)$ to be the number of powers in that
sum, and set $s_2(0)=0$. For example,
\[
 13=8+4+1=(1101)_2,\qquad s_2(13)=3.
\]
Two properties will be used:
\begin{equation}\label{eq:digits}
 s_2(2r)=s_2(r),\qquad
 s_2(r_1+\cdots+r_k)\le\sum_{i=1}^{k}s_2(r_i)
 \quad(r_i\ge0).
\end{equation}
Multiplication by two shifts every binary digit one place, proving the
first identity. For the inequality, expand each summand into powers of
two. Whenever a power occurs twice, replace the two copies by one copy of
the next power. Each replacement preserves the value and decreases the
number of summands. The process terminates at the binary expansion of the
total, which proves the inequality.

\subsection{The divisibility theorem}

Let $B_m=x+\sum_{n\ge2}b_n(m)x^n$ be the compositional inverse of
$A_m$, supplied by Lemma~\ref{lem:inverse}, and set $b_1(m)=1$.

\begin{theorem}\label{thm:binary}
For every integer $m\ge2$ and every $n\ge1$,
\begin{equation}\label{eq:binary}
 m^{s_2(n-1)}\mid a_n(m)
 \qquad\text{and}\qquad
 m^{s_2(n-1)}\mid b_n(m).
\end{equation}
\end{theorem}

\begin{proof}
Fix $m\ge2$. Remove the initial factor of $x$ from the two inverse series:
\[
 A_m(x)=xF(x),\qquad B_m(x)=xG(x),
\]
\[
 F(x)=\sum_{r\ge0}f_rx^r,\qquad
 G(x)=\sum_{r\ge0}g_rx^r,\qquad f_0=g_0=1.
\]
Then $f_r=a_{r+1}(m)$ and $g_r=b_{r+1}(m)$, so the shifted index $r$
is the integer whose binary digits occur in the desired bound.

Apply $B_m$ to \eqref{eq:general}. Using $B_m(A_m(x))=x$ gives
\[
 xA_m-mxA_m^2=B_m(x^2).
\]
Hanna's entries record this identity for $m=1,2,3,4,5$
\cite{oeis1,oeis2,oeis3,oeis4,oeis5}. We use it here for general $m$.
Substitution of $A_m=xF$ and $B_m=xG$, followed by cancellation of the
common factor $x^2$, yields
\begin{equation}\label{eq:FG}
 F(x)=G(x^2)+mxF(x)^2.
\end{equation}
Cancellation here simply shifts every coefficient two places. For $r\ge1$,
comparison of coefficients gives
\begin{equation}\label{eq:frec}
 f_r=\begin{cases}g_{r/2},&r\text{ even},\\0,&r\text{ odd}\end{cases}
       +m\sum_{i+j=r-1}f_if_j,
\end{equation}
where $i,j$ are nonnegative integers.

The other inverse identity, $A_m(B_m(x))=x$, gives
\[
 1=\sum_{j\ge0}f_jx^jG(x)^{j+1}.
\]
In degree $r\ge1$, the $j=0$ term contributes $g_r$ and the $j=r$
term contributes $f_r$, because $G(0)=1$. Thus
\begin{equation}\label{eq:grec}
 g_r=-f_r-\sum_{j=1}^{r-1}f_j[x^{r-j}]G(x)^{j+1}.
\end{equation}

We prove by induction on $r$ that
\begin{equation}\label{eq:simultaneous}
 m^{s_2(r)}\mid f_r\quad\text{and}\quad m^{s_2(r)}\mid g_r.
\end{equation}
The case $r=0$ says that $1$ divides $f_0=g_0=1$.
Suppose the assertions hold at every smaller index.

First consider $f_r$ in \eqref{eq:frec}. If $r$ is even, then
$r/2<r$ and $s_2(r/2)=s_2(r)$, so the term $g_{r/2}$ has the required
divisor. Each product $mf_if_j$ uses $i,j<r$ and is divisible by
$m^{1+s_2(i)+s_2(j)}$. The digit inequality \eqref{eq:digits} gives
\[
 1+s_2(i)+s_2(j)\ge s_2(1+i+j)=s_2(r).
\]
Every term in \eqref{eq:frec} is therefore divisible by $m^{s_2(r)}$,
and so is their sum $f_r$.

Now consider $g_r$ in \eqref{eq:grec}. The first term $-f_r$ has just
been handled. For $1\le j<r$, expansion of $G^{j+1}$ expresses each
remaining summand as a sum of products
\[
 f_jg_{i_1}\cdots g_{i_{j+1}},\qquad
 j+i_1+\cdots+i_{j+1}=r,
 \qquad i_\ell\ge0.
\]
Each $g$ index is at most $r-j<r$. By the induction hypothesis the product
is divisible by
\[
 m^{s_2(j)+s_2(i_1)+\cdots+s_2(i_{j+1})},
\]
and its exponent is at least $s_2(r)$ by \eqref{eq:digits}.
Taking sums and negatives preserves divisibility, proving the assertion
for $g_r$. This completes the induction: at each stage we establish $f_r$
first and then $g_r$.

Finally, replace $r$ by $n-1$ in \eqref{eq:simultaneous}.
\end{proof}

\begin{remark}\label{rem:polynomial}
The same argument works before the parameter is assigned an integer value.
Let $t$ be an indeterminate and construct $A_t$ and $B_t$ over $\Z[t]$.
All coefficient recurrences above use only sums, products, and negatives. The induction
therefore proves that both $a_n(t)$ and $b_n(t)$ are multiples of
$t^{s_2(n-1)}$ in $\Z[t]$. In other words, neither polynomial has a
nonzero term of degree below $s_2(n-1)$ in the parameter $t$.
\end{remark}

For illustration, the original parameter $m=3$ gives the following bounds.
\begin{center}
\begin{tabular}{@{}rrrr@{}}
\toprule
$n$ & Binary expansion of $n-1$ & $s_2(n-1)$ & Guaranteed divisor of $a_n(3)$\\
\midrule
$2$ & $1$ & $1$ & $3$\\
$8$ & $111$ & $3$ & $27$\\
$9$ & $1000$ & $1$ & $3$\\
$32$ & $11111$ & $5$ & $243$\\
$64$ & $111111$ & $6$ & $729$\\
\bottomrule
\end{tabular}
\end{center}
In general, $2^k-1$ has $k$ binary ones, so
$m^k\mid a_{2^k}(m)$ for every $k\ge0$. The table also shows why the
guaranteed divisor need not grow with the index: a power of two has only
one binary one.

\subsection{Exactly where the smallest prime exponent occurs}

For a prime $p$ and a nonzero integer $q$, write $v_p(q)$ for the number
of factors of $p$ in $q$. Equivalently, it is the largest integer $e\ge0$
such that $p^e\mid q$. For example, $v_3(45)=2$. Set $v_p(0)=+\infty$.
This notation lets us distinguish a lower bound from an exact exponent.

\begin{corollary}\label{cor:valuation}
Let $m\ge2$, let $p$ be a prime dividing $m$, and put $e=v_p(m)$. Then
\[
 v_p(a_n(m))\ge e\,s_2(n-1)\qquad(n\ge1).
\]
For $n\ge2$, equality with the smallest possible value $e$ is characterized by
\begin{equation}\label{eq:minval}
 v_p(a_n(m))=e
 \quad\Longleftrightarrow\quad
 n=2^j+1\text{ for some }j\ge0.
\end{equation}
\end{corollary}

\begin{proof}
The lower bound follows immediately from Theorem~\ref{thm:binary}.
At $n=2^j+1$, Theorem~\ref{thm:modsquare} gives
\[
 a_n(m)=m\bigl((-1)^j+mq\bigr)
\]
for some integer $q$. The factor in parentheses is $1$ or $-1$ modulo $p$,
so it is not divisible by $p$. Thus $a_n(m)$ has exactly $e$ factors of $p$.
At any other index $n\ge2$, the positive integer $n-1$ is not a power
of two and hence has at least two binary ones. The lower bound is then
at least $2e>e$.
\end{proof}


\section{Specialization to A396844}
Substitute $m=4$ in Theorems~\ref{thm:parity},
\ref{thm:modsquare}, and \ref{thm:binary}, and in
Corollary~\ref{cor:valuation}. In the residue formula, write
$j=2k$ or $j=2k+1$. Then $2^j=4^k$ or $2\cdot4^k$ respectively,
and $(-1)^j=1$ or $-1$. This proves precisely the two exceptional
classes stated at the start, with every other index handled by the
second case of Theorem~\ref{thm:modsquare}. The powers of two are
strictly increasing, so these exceptional classes are disjoint.

\section{Formal proof and scope}
The accompanying entry file is \path{lean/Entries/A396844.lean}.
It states the results with explicit Lean types and supplies proof terms
using the complete shared implementation included in the collection.
The manuscript is self-contained mathematically; the Lean entry file uses
those ordinary library imports rather than duplicating the full foundation.
The complete original acceptance layer is retained in \path{lean/Hanna/Acceptance.lean}. The shared endpoints \path{Hanna.Acceptance.inverse_binary_divisibility_integer}
and \path{Hanna.Acceptance.minimum_prime_exponent_iff} give the general
inverse-divisibility and minimum-exponent assertions used in this paper.
The counterexample at $n=17$ follows by taking $k=2$ in the checked
\path{residues_parameter_four} classification.

The entry's named theorem interfaces are listed below. The namespace of
each is \path{FinalCollection.A396844}.
\begin{itemize}[leftmargin=1.5em,itemsep=1pt]
\item \path{source_existence_uniqueness}
\item \path{even_positive_tail}
\item \path{binary_divisibility}
\item \path{residues_parameter_four}
\end{itemize}
The preserved dependencies also contain all supporting constructions,
coefficient identities, map-compatibility results, and inverse proofs.
The collection pins Lean 4.19.0 and Mathlib revision
\path{c44e0c8ee63ca166450922a373c7409c5d26b00b}.
The verification report records fresh compilation and the transitive axiom
audit. Finite numerical checks corroborate the definitions and indexing;
the universally quantified Lean theorems establish the all-index claims.
No assertion of worldwide novelty, OEIS acceptance, publication, or completed
human review follows from these checks. The prose may use elementary
identities that the formal proof establishes by a different algebraic route.
\begin{thebibliography}{9}\setlength{\itemsep}{2pt}\setlength{\parsep}{0pt}
\bibitem{oeis1} P. D. Hanna, \emph{A265940}, The On-Line Encyclopedia of
Integer Sequences, 19 December 2015. \url{https://oeis.org/A265940}.
Accessed 22 October 2026.
\bibitem{oeis2} P. D. Hanna, \emph{A271958}, The On-Line Encyclopedia of
Integer Sequences, 17 April 2016. \url{https://oeis.org/A271958}.
Accessed 22 October 2026.
\bibitem{oeis3} P. D. Hanna, \emph{A396843}, The On-Line Encyclopedia of
Integer Sequences, 1 July 2026. \url{https://oeis.org/A396843}.
Accessed 22 October 2026.
\bibitem{oeis4} P. D. Hanna, \emph{A396844}, The On-Line Encyclopedia of
Integer Sequences, 3 July 2026. \url{https://oeis.org/A396844}.
Accessed 22 October 2026.
\bibitem{oeis5} P. D. Hanna, \emph{A396845}, The On-Line Encyclopedia of
Integer Sequences, 3 July 2026. \url{https://oeis.org/A396845}.
Accessed 22 October 2026.
\end{thebibliography}
\end{document}
